\subsection{Collective Coupling and Operational Geometry}
  \label{subsec:collective-coupling-operational-geometry}

  Curvature in a relational framework arises as a collective effect.
  Geometric properties emerge from how relational variations influence one
  another across extended regions, without postulating a fundamental
  gravitational field.

  In approximately homogeneous regimes, variations propagate uniformly.
  Localized irregularities modify this collective response,
  which can be summarized by an effective coupling characterizing
  the stiffness of the relational structure under relative deformations.

  Distance is defined operationally:
  regions are proximate if relational variations can be efficiently correlated.
  In weakly inhomogeneous continuum regimes,
  this operational proximity admits a compact geometric representation.
  An effective metric summarizes the collective response of the system,
  without constituting an independent degree of freedom.

  Curvature thus appears as a macroscopic descriptor of how localized
  relational structure modulates global connectivity.

  \paragraph{Spectral convergence and continuum limit.}

    Under standard regularity assumptions
    (uniform density, bounded degree growth, approximate isotropy),
    the graph Laplacian $\Delta_G$ converges,
    in the strong resolvent sense,
    to a Laplace--Beltrami operator $\Delta_{\mathcal M}$
    on an effective manifold $\mathcal M$.

    Operationally, low-lying eigenmodes of $\Delta_G$
    approximate those of $\Delta_{\mathcal M}$
    below a spectral cutoff $\lambda_*$.
    Beyond this scale, smooth geometry ceases to apply.
    Such convergence results are well established in spectral graph theory
    and justify the continuum interpretation of the admissible spectral sector.

  \paragraph{Spectral dimension as emergent observable.}

    Effective dimensionality is inferred from spectral data.
    Using the heat kernel
    \begin{equation}
      K(t) = \sum_n e^{-t \lambda_n},
    \end{equation}
    the spectral dimension is defined by
    \begin{equation}
      d_s(t) = -2 \frac{d \log K(t)}{d \log t}.
    \end{equation}

    In admissible regimes, $d_s(t)$ exhibits a stable large-scale plateau.
    Configurations relevant to the gravitational regime display
    $d_s \simeq 4$ robustly across discretizations,
    without imposing dimensionality \emph{a priori}.

  \paragraph{Role of the spectral scale $\lambda_*$.}

    The cutoff $\lambda_*$ defines the upper limit of spectral admissibility.
    Above this scale, spectral locality breaks down and smooth geometric
    reconstruction becomes ill-defined.

    This ultraviolet breakdown signals the transition to a
    pre-geometric regime rather than a failure of the framework.
